% ================= COPERTINA =================
\thispagestyle{empty}
\begin{tikzpicture}[remember picture,overlay]
  \fill[navy] (current page.north west) rectangle ([yshift=-15.5cm]current page.north east);
  \fill[teal] ([yshift=-15.5cm]current page.north west) rectangle ([yshift=-16.1cm]current page.north east);
  % gaussiane decorative
  \begin{scope}[shift={([xshift=-1cm,yshift=-15.5cm]current page.north west)},x=1cm,y=1cm]
    \fill[white,opacity=.07,domain=0:23,samples=120] (0,0) -- plot (\x,{9*exp(-((\x-8)^2)/(2*2.2^2))}) -- (23,0) -- cycle;
    \fill[teal,opacity=.30,domain=0:23,samples=120] (0,0) -- plot (\x,{6.5*exp(-((\x-14.5)^2)/(2*2.6^2))}) -- (23,0) -- cycle;
    \draw[white,opacity=.55,line width=1.2pt,domain=0:23,samples=120] plot (\x,{9*exp(-((\x-8)^2)/(2*2.2^2))});
    \draw[ambra,opacity=.9,line width=1.2pt,domain=0:23,samples=120] plot (\x,{6.5*exp(-((\x-14.5)^2)/(2*2.6^2))});
    \draw[white,dashed,opacity=.6] (11.55,0) -- (11.55,8.6);
    \node[white,opacity=.8,font=\headfont\small] at (11.55,8.9) {$x_B$};
  \end{scope}
  \node[anchor=north west,font=\headfont\small\bfseries,text=ambra] at ([xshift=2.3cm,yshift=-2.2cm]current page.north west) {MACHINE LEARNING \,·\, ELEMENTI DI INTELLIGENZA ARTIFICIALE};
  \node[anchor=north west,font=\headblack\fontsize{46}{50}\selectfont,text=white,align=left] at ([xshift=2.2cm,yshift=-3.0cm]current page.north west) {Classificatore\\[2pt]di Bayes};
  \node[anchor=north west,font=\headlight\Large,text=white!85!teal,align=left] at ([xshift=2.3cm,yshift=-6.9cm]current page.north west) {Probabilità, decisione ottima e approccio parametrico};
  % formula in copertina
  \node[anchor=north east,text=white,font=\Large] at ([xshift=-2.3cm,yshift=-17.2cm]current page.north east) {};
\end{tikzpicture}

\vspace*{15.2cm}
\begin{center}
\begin{tcolorbox}[enhanced,width=.88\linewidth,colback=white,colframe=navy!35,arc=6pt,boxrule=.8pt,halign=flush center,top=10pt,bottom=10pt]
\large$\displaystyle P(w_i\,|\,\x) = \frac{p(\x\,|\,w_i)\,P(w_i)}{p(\x)}\qquad\quad b=\argmax_{i=1,\dots,s} P(w_i\,|\,\x)$
\end{tcolorbox}
\vspace{1.0cm}

{\headfont\large\color{navy}\bfseries Dispensa di studio\par}
\vspace{2pt}
{\headfont\color{grigio}Appunti ragionati sulle slide del corso \,·\, A.A. 2026–2027\par}
\vspace{1.2cm}

\begin{tikzpicture}
  \foreach \c/\t [count=\i] in {teal/Intuizione,navy/Tecnica,ambra/Esempi,rosso/Attenzione,viola/ML}{
    \node[fill=\c,text=white,font=\headfont\bfseries\footnotesize,rounded corners=3pt,minimum width=2.3cm,minimum height=.6cm] at (\i*2.65,0) {\t};
  }
\end{tikzpicture}
\end{center}

% ================= COME USARE =================
\clearpage
\thispagestyle{plain}
{\headfont\bfseries\Huge\color{navy}Come usare questa dispensa\par}
\vspace{4pt}
{\color{teal}\rule{3cm}{2.5pt}}
\vspace{8pt}

La dispensa segue le 62 slide della lezione \emph{Classificatore di Bayes} divise in \hl{13 sezioni}. Ogni sezione si apre con il riferimento alle slide e spiega il contenuto due volte: prima in modo \key{intuitivo}, poi in modo \key{tecnico}, con formule, esempi numerici verificati e grafici. Alla fine di ogni sezione trovi un riquadro grigio di \hl{appunti brevi}, organizzato per numero di slide, da copiare accanto alle slide originali.

I riquadri colorati hanno sempre lo stesso significato:

\begin{intuizione}
L'idea di fondo, senza formule: analogie, immagini mentali, "perché funziona".
\end{intuizione}
\begin{tecnica}
Definizioni precise, formule e passaggi matematici.
\end{tecnica}
\begin{esempio}
Calcoli svolti passo per passo. Tutti i numeri sono stati ricontrollati.
\end{esempio}
\begin{attenzione}
Errori tipici, trappole di notazione e precisazioni sulle slide.
\end{attenzione}
\begin{mlbox}
Dove ritroverai il concetto nel Machine Learning moderno (k-NN, PCA, LDA, regolarizzazione\dots).
\end{mlbox}
\begin{formulat}{Formula chiave}
I riquadri bianchi bordati contengono le formule da sapere a memoria.
\end{formulat}

\vspace{4pt}
{\headfont\bfseries\color{navy}Notazione.}\;
Vettori in grassetto ($\x$, $\mmu$), matrici in grassetto maiuscolo ($\SSigma$). $\x_i$ è l'$i$-esimo pattern del training set; $x^j$ è la $j$-esima componente (feature). L'apice $\mathsf{t}$ indica la trasposta. $P(\cdot)$ maiuscola = probabilità, $p(\cdot)$ minuscola = densità.

\clearpage
\thispagestyle{plain}
\tableofcontents
